\chapter{Executable quantum codes and neural circuits on dessins}
\label{ch:quantum-surfaces}

\section{From a drawing to a quantum computation}

A dessin becomes relevant to a quantum circuit only after the underlying
combinatorics determine concrete Hilbert spaces, operators, and gates.  This
chapter carries out that construction.  It supplies oriented cellulations,
their clean-dessin permutations, binary boundary matrices, commuting Pauli
checks, an executable encoder, and a small trainable circuit whose output
can be independently derived.  The smallest live example has eighteen
data qubits, rather than a picture of eighteen unrelated two-level systems.
Its state has entanglement enforced by the same incidence relations that
describe the torus.

The construction is the homological surface-code construction introduced
by Kitaev, within the Calderbank--Shor--Steane framework
\cite{QSKitaev2003,QSCalderbankShor1996,QSSteane1996}.  Here the contribution
is a reproducible connection to the manuscript's orders, complexes and
dessins, together with exact examples and verification artifacts.  Existing
hypermap-code literature already studies broader embedded combinatorial
objects; the mere use of a dessin does not establish a new code family or
better parameters \cite{QSLeslie2013,QSSarvepalli2014}.

There are three computational representations.  The browser executes the
eighteen-qubit gate sequence as a full complex state vector.  Larger
Clifford computations can use stabilizer or Pauli-frame representations.
GPU state-vector experiments use the resources stated in their recorded
reports.  These representations answer different computational questions.
The size of a stabilizer simulation is not evidence that the same machine
stored an arbitrary state on that many qubits.

\section{Oriented cells and clean-dessin permutations}

Let $\Sigma$ be a connected, closed, oriented surface with a finite cellular
embedding of a graph.  Let $V,E,F$ denote its numbers of vertices, edges
and disk faces.  An oriented edge has a tail and a head; reversing it
produces the other dart of the same unoriented edge.  There are therefore
$2E$ darts.  Put a black vertex at every original vertex and a white vertex
at the midpoint of every edge.  The resulting bipartite graph is clean:
every white vertex has degree two.

Write $\sigma_1$ for dart reversal and $\varphi$ for the permutation that
follows the oriented boundary of a face.  Our convention is
\[
 \sigma_0=\varphi\sigma_1,\qquad
 \sigma_\infty=(\sigma_0\sigma_1)^{-1}=\varphi^{-1},
\]
with products acting from right to left.  The cycles of $\sigma_0$ are
the rotations at black vertices, the transpositions of $\sigma_1$ are
white vertices, and the cycles of $\sigma_\infty$ are faces with reversed
boundary orientation.  The code checks these identities as permutations;
it does not infer them from the apparent shape of a rendered mesh.

\begin{proposition}[Genus from permutation cycles]
For the connected oriented cellulation above,
\[
 2-2g=c(\sigma_0)+c(\sigma_1)+c(\sigma_\infty)-2E
      =V-E+F,
\]
where $c$ denotes the number of cycles.
\end{proposition}
\begin{proof}
Subdividing each edge inserts one white vertex and replaces that edge by
two bipartite edges.  The new numbers of vertices and edges are $V+E$ and
$2E$; the disk faces are unchanged.  Its Euler characteristic is therefore
$(V+E)-2E+F=V-E+F$.  The three cycle counts are $V,E,F$, by their
definitions.  The classification formula for a connected closed oriented
surface gives the final value $2-2g$.
\end{proof}

An abstract rotation system needs a manifold check.  The boundary cycles
must use every dart exactly once, and all darts based at a given vertex
must form one rotation cycle.  Multiple rotation cycles at one declared
vertex would create a nonmanifold identification.  The implementation
checks these conditions, checks graph connectivity, and requires every
unoriented edge to occur twice in face boundaries with opposite signs.
Thus the Euler characteristic is attached to an actual closed oriented
cellulation, not merely to three integers satisfying an equation.

The Riemann existence theorem realizes a transitive permutation triple
with product one as monodromy of a branched covering of the Riemann sphere.
For a dessin, its branch values can be chosen as $0,1,\infty$.  Belyi's
theorem relates this situation to curves defined over algebraic numbers
\cite{QSBelyi1980}.  These are external existence theorems.  The generated
permutations are explicit; algebraic coefficients of the higher-genus
Belyi maps are not computed by this implementation.  Nor does a Euclidean
rendering specify the uniformizing metric of that Belyi pair.

\section{An explicit elliptic Belyi model}

The algebraic curve
\[
 E_0:\quad y^2=x^3-x
\]
provides a direct connection with elliptic-curve theory.  Its smooth
projective completion has genus one, and the meromorphic function
$\beta=x^2$ has degree four.

\begin{proposition}[A clean degree-four elliptic dessin]
The branch values of $\beta:E_0\to\mathbb P^1$ lie in
$\{0,1,\infty\}$.  Its ramification partitions over these values are
$(4)$, $(2,2)$ and $(4)$, respectively.  Its dessin has one black vertex,
two white vertices and one face.
\end{proposition}
\begin{proof}
At $(0,0)$, the curve equation gives $x=-y^2+O(y^6)$, so $\beta$ has
a zero of order four.  At $(r,0)$ for $r=1$ or $r=-1$, the derivative of
$x^3-x$ is two.  Thus $x-r=y^2/2+O(y^4)$, and
$\beta-1=2r(x-r)+O((x-r)^2)$ has order two.  At the point at infinity,
$x$ has a pole of order two, so $\beta$ has a pole of order four.
These fibers already account for degree four.  Away from them, a critical
point would require $x=0$ or a critical point of the degree-two projection
to the $x$-line; the latter occur at $x=0,1,-1,\infty$, already listed.
The multiplicities give the claimed partitions and cycle counts.
\end{proof}

Topologically the dessin is a two-loop graph on a torus with disk
complement.  Choose its two loop classes as generators of the fundamental
group.  The connected unramified cover associated with
$L_x\mathbb Z\oplus L_y\mathbb Z\subseteq\mathbb Z^2$ has degree
$L_xL_y$.  Pulling back the dessin gives the square torus cellulation
used below.  The covering becomes holomorphic for the pulled-back complex
structure, so composing it with $\beta$ gives a clean Belyi map of degree
$4L_xL_y$.  This supplies a genuine elliptic-surface interpretation of
the torus family without requiring numerical evaluation of elliptic
functions during circuit execution.

For comparison, an explicit higher-genus family is
\[
 C_g:\quad y^2=1+x^{2g+1},\qquad \beta_g=-x^{2g+1}=1-y^2,
 \quad g\geq1.
\]
The double cover of the $x$-sphere has $2g+2$ simple branch points,
including infinity, so its smooth completion has genus $g$.  The displayed
Belyi map has two black vertices of valency $2g+1$, $2g+1$ white vertices
of valency two, and one face.  Locally this follows from the two points
over $x=0$, the simple roots of $1+x^{2g+1}$, and the pole orders at
infinity.  Its degree is $2(2g+1)$ and its Euler characteristic is
$2-(2g+1)+1=2-2g$.

This last dessin is an instructive warning.  Collapsing its white
subdivision produces two vertices, $2g+1$ parallel edges and one face.
The face boundary is zero modulo two, so the corresponding code has no
nontrivial $Z$ checks and distance one.  Genus alone does not make an
error-correcting code useful.  The executable higher-genus cases instead
use the nondegenerate connected-sum cellulations constructed below.

\section{The cellular chain complex is the stabilizer specification}

Work over $\mathbb F_2$.  Give the cellular vector spaces their cell bases:
\[
 C_2=\mathbb F_2^F\xrightarrow{\partial_2}
 C_1=\mathbb F_2^E\xrightarrow{\partial_1}\mathbb F_2^V=C_0.
\]
The boundary of an edge is the sum of its endpoints; a loop contributes
twice and hence zero.  The boundary of a face is the parity of its edge
occurrences.  Since the boundary of a closed polygon has each vertex an
even number of times, $\partial_1\partial_2=0$.  Define
\[
 H_X=\partial_1,\qquad H_Z=\partial_2^{\mathsf T}.
\]
Their rows are the vertex stars and face boundaries.  In the actual files
the rows are sparse lists of qubit indices.  They are not dense floating
matrices subsequently rounded to binary values.

Place one qubit on each original edge.  For $a,b\in\mathbb F_2^E$, put
\[
 X(a)=\prod_eX_e^{a_e},\qquad Z(b)=\prod_eZ_e^{b_e}.
\]
The single-qubit identity $XZ=-ZX$ gives
\[
 X(a)Z(b)=(-1)^{a\cdot b}Z(b)X(a).
\]
Let $A_v$ be the $X$ product on the star of vertex $v$, and $B_f$ the
$Z$ product on the boundary of face $f$.

\begin{theorem}[Commutation from the chain identity]
All $A_v$ and $B_f$ commute.  Their common $+1$ eigenspace is a CSS
code with
\[
 k=E-\operatorname{rank}H_X-\operatorname{rank}H_Z
\]
logical qubits.
\end{theorem}
\begin{proof}
Two products of $X$ operators commute, as do two products of $Z$
operators.  The commutation sign of $A_v$ with $B_f$ is the $(v,f)$
entry of $H_XH_Z^{\mathsf T}=\partial_1\partial_2=0$.
Choose independent rows in the two matrices.  Their Pauli generators
are independent jointly: an identity product must have zero $X$ mask
and zero $Z$ mask separately.  Ordering all $X$ factors before all $Z$
factors introduces no sign because these groups commute.  Thus the
generated stabilizer group contains $2^r$ elements, where
$r=\operatorname{rank}H_X+\operatorname{rank}H_Z$, and does not contain
$-I$.

The projector onto the common $+1$ eigenspace is
$P=2^{-r}\sum_{S\in\mathcal S}S$.  Multiplication in the finite
group proves $P^2=P$, and Hermiticity follows because its elements
are Hermitian involutions.  Every nonidentity Pauli has trace zero,
while $I$ has trace $2^E$.  Therefore
$\dim\operatorname{im}P=\operatorname{tr}P=2^{E-r}=2^k$.
\end{proof}

\begin{theorem}[Logical qubits on a closed oriented surface]
For a connected closed oriented surface of genus $g$,
\[
 \operatorname{rank}H_X=V-1,\qquad
 \operatorname{rank}H_Z=F-1,\qquad k=2g.
\]
\end{theorem}
\begin{proof}
A vector in $\ker\partial_1^{\mathsf T}$ assigns equal binary
values to the two endpoints of every edge.  Connectivity forces it
to be constant.  Its kernel has dimension one, so
$\operatorname{rank}\partial_1=V-1$.

Assign a coefficient $a_f$ to each face.  In
$\partial_2\sum_fa_ff$, the coefficient of an edge separating faces
$f$ and $f'$ is $a_f+a_{f'}$.  The dual graph is connected: a path
between face interiors can be perturbed to cross edges transversely
and avoid vertices.  Consequently the kernel consists exactly of
the constant face assignments, and
$\operatorname{rank}\partial_2=F-1$.  The all-face sum really has
zero boundary because each edge occurs twice.  Substitution in the
CSS formula gives $k=E-V-F+2=2-(V-E+F)=2g$.
\end{proof}

This is a quantitative use of topology, but its quantity is the dimension
of a finite code space.  It does not compute a classical cryptographic
secret or a continuum field-theory mass.  The theorem also explains why
one must retain the dependent checks: the products of all star checks
and of all face checks are each identity, although the checks remain
useful local observables.

\section{Logical operators, distance and syndrome}

A $Z$ string commutes with every star precisely when its edge mask belongs
to $\ker\partial_1$.  Multiplying it by a face stabilizer adds a boundary.
Thus its logical class lies in
\[
 H_1(\Sigma;\mathbb F_2)=\ker\partial_1/\operatorname{im}\partial_2.
\]
Similarly $X$ strings are classified by
$\ker H_Z/\operatorname{row}H_X$, the corresponding cohomological
quotient.  The dot product is their commutation pairing.  A nonzero pairing
means that the associated logical $X$ and $Z$ anticommute.  This provides
the operator interpretation of the $2g$ homology generators.

For an error $X(a)Z(b)$, its syndrome is
\[
 s_X=H_Xb,\qquad s_Z=H_Za.
\]
Each syndrome bit records an anticommuting check.  A single $Z$ error
creates star defects at the edge endpoints, and a single $X$ error creates
defects on its two adjacent faces.  A recovery with the same syndrome
can still differ from the error by a nontrivial logical loop.  Therefore
displaying an error syndrome is not by itself a demonstration of a decoder.
Repeated noisy syndrome extraction and decoding require the additional
protocol assumptions studied in surface-code memory analyses
\cite{QSDennis2002}.

For the $L_x\times L_y$ square torus with $L_x,L_y\geq3$, index the
horizontal and vertical edges by
\[
 h(x,y)=yL_x+x,\qquad v(x,y)=L_xL_y+yL_x+x,
\]
using coordinates modulo the two lengths.  A face at $(x,y)$ has oriented
boundary
\[
 h(x,y)+v(x+1,y)-h(x,y+1)-v(x,y).
\]
There are $L_xL_y$ vertices and faces, and $2L_xL_y$ qubits.

\begin{proposition}[Distance of the rectangular torus code]
The square torus construction has parameters
\[
 [[2L_xL_y,\ 2,\ \min(L_x,L_y)]].
\]
\end{proposition}
\begin{proof}
The logical dimension is the preceding theorem.  A binary cycle is a
union of closed edge paths.  If its homology class is nonzero, at least
one constituent path has odd winding in a periodic direction.  Lift that
path to the square lattice in the plane.  Its endpoint displacement has
a nonzero coordinate equal to an integral multiple of $L_x$ or $L_y$.
Each edge changes one coordinate by one, so the path has length at least
$\min(L_x,L_y)$.  The same argument applies to the dual square lattice.
Straight horizontal and vertical loops attain the corresponding lengths,
so the minimum is attained.

For completeness, a mixed Pauli logical operator cannot improve this
minimum.  Its $X$ and $Z$ masks satisfy their separate centralizer
conditions.  If both were stabilizer masks, their product would also
be a stabilizer.  At least one is a nontrivial logical mask, and the
weight of their union is at least the weight of that component.
\end{proof}

On the eighteen-qubit torus, examples of conjugate logical pairs are
\[
 Z_{\rm h}=Z_0Z_1Z_2,\quad X_{\rm h}=X_0X_3X_6,
 \qquad
 Z_{\rm v}=Z_9Z_{12}Z_{15},\quad X_{\rm v}=X_9X_{10}X_{11}.
\]
The displayed matched pairs overlap once; unmatched pairs have zero
overlap.  All commute with the stabilizers.  The backend also checks
distance independently by exhaustively testing increasing Pauli-support
weights modulo the stabilizer row spaces.  For CSS codes it suffices to
test the two pure types, by the last paragraph of the proof.

\section{Higher genus by a verified connected sum}

To obtain genus two, take two copies of the $3\times3$ torus and remove
one open square face from each.  If their boundary vertices are
$a_0,\ldots,a_3$ and $b_0,\ldots,b_3$ in their original face
orientations, identify $a_i$ with $b_{-i}$, with indices modulo four.
This reverses the boundary orientation.  The boundary edge
$a_i a_{i+1}$ is paired with the reverse of
$b_{-i-1}b_{-i}$.

The new counts are
\[
 V=9+9-4=14,\quad E=18+18-4=32,\quad F=9+9-2=16.
\]
Thus $\chi=-2$ and $g=2$.  Pairing edges explicitly is important:
identifying every edge with the same unordered endpoint pair would
incorrectly merge distinct parallel edges in more general maps.  The
implementation retains edge identities, reverses their incidence signs
when required, and reconstructs every face and vertex rotation cycle.

A second connected sum with a third torus gives
$(V,E,F)=(19,46,23)$ and genus three.  The selected remaining face has
four distinct vertices, so its closure is an embedded disk even when it
shares vertices with the earlier seam.  The same manifold checks are
repeated after gluing.  The counts are not accepted as a substitute for
these checks.

\begin{longtable}{lrrrrrr}
\toprule
Cellulation & $E$ & $g$ & $\operatorname{rank}H_X$ &
$\operatorname{rank}H_Z$ & $k$ & $d$\\\midrule
$3\times3$ torus &18&1&8&8&2&3\\
$3\times5$ torus &30&1&14&14&2&3\\
Two tori, connected sum &32&2&13&15&4&3\\
Three tori, connected sum &46&3&18&22&6&3\\
$32\times32$ torus &2048&1&1023&1023&2&32\\\bottomrule
\end{longtable}

The first four distances are checked by exhaustive search through weight
three, including proof by enumeration that smaller nontrivial logical
supports do not exist.  The last distance uses the rectangular-torus
theorem, not a search over $2^{2048}$ masks.  At a genus-two seam, some
star checks have weight six; the genus-three example has star weights up
to eight.  All their face checks have weight four.  These constructions
increase the number of logical qubits, but the chosen small cells retain
distance three.  They do not establish an improved asymptotic rate-distance
tradeoff.

\section{A deterministic encoder with a complete proof}

Let $W=\operatorname{row}H_X\subseteq\mathbb F_2^E$ and let
$r=\dim W$.  The prepared code state is
\begin{equation}
 |\Omega\rangle=2^{-r/2}\sum_{w\in W}|w\rangle.
 \label{eq:surface-encoded-state}
\end{equation}
Distinct binary strings label orthonormal computational basis vectors,
so the normalization follows from $|W|=2^r$.

\begin{proposition}[Stabilizers and logical sector of the prepared state]
Every star and face stabilizes $|\Omega\rangle$.  Every $Z(b)$ with
$b\in\ker H_X$ also has eigenvalue $+1$ on this state.  Thus the state
chooses the simultaneous $+1$ sector of a full set of logical $Z$
operators.
\end{proposition}
\begin{proof}
An $X$ star adds an element of $W$ to each summation label, permuting
$W$ and leaving its uniform sum invariant.  For $b\in\ker H_X$,
every $w\in W$ satisfies $b\cdot w=0$, so
$Z(b)|w\rangle=(-1)^{b\cdot w}|w\rangle=|w\rangle$.
The face masks lie in $\ker H_X$ by the chain identity.  The remaining
homologically nontrivial masks provide the logical $Z$ conditions.
\end{proof}

Reduce $H_X$ to a binary reduced row-echelon basis
$r_1,\ldots,r_r$, with pivot positions $p_1,\ldots,p_r$.
On an all-zero register, apply a Hadamard to each pivot.  Then, for every
nonpivot position $q$ with $(r_i)_q=1$, apply
$\mathrm{CX}_{p_i\to q}$.  This is the actual gate sequence in the JSON
files, not merely a reference to an abstract encoding unitary.

\begin{theorem}[Correctness of the encoder]
The specified Hadamard--CNOT circuit prepares
\eqref{eq:surface-encoded-state} exactly in the ideal gate model.
\end{theorem}
\begin{proof}
After the Hadamards, the pivot bits range uniformly over
$u=(u_1,\ldots,u_r)\in\mathbb F_2^r$, with amplitude $2^{-r/2}$,
and all other bits are zero.  Reduced row-echelon form makes the pivot
submatrix the identity.  No later CNOT targets a pivot.  Therefore its
control retains $u_i$, and the final nonpivot bit $q$ is
$\sum_i u_i(r_i)_q$.  The entire output string is
$\sum_i u_ir_i$.  Independence of the rows makes this a bijection
from $\mathbb F_2^r$ to $W$.  CNOTs permute basis states without phases,
which proves the claimed uniform positive sum.
\end{proof}

The circuits have respectively $42,74,82,134$ gates for the eighteen,
thirty, thirty-two and forty-six-qubit cases.  The eighteen-qubit circuit
uses eight Hadamards and thirty-four CNOTs.  This is an algebraic encoder;
some CNOT endpoints are distant in the surface drawing.  A nearest-neighbor
hardware compilation would need routing, scheduling and extra gates.
Local stabilizer supports do not automatically make this particular
preparation circuit local.

The implementation uses little-endian basis indices: qubit $q$ is bit $q$
of the array index.  Each gate specifies its name and target, a control for
CNOT, and an angle for rotations.  This convention is shared by the
browser and the CPU/GPU interfaces.  An inconsistent bit convention could
produce a normalized state while assigning its amplitudes to the wrong
qubits, so normalization alone is not an adequate validation.

\section{A finite Hamiltonian with a provable gap}

The same checks define the positive Hamiltonian
\[
 H_{\rm code}=\sum_v\frac{I-A_v}{2}
              +\sum_f\frac{I-B_f}{2}.
\]
Every term is an orthogonal projector and all terms commute.  The ground
space is exactly the CSS code space.

\begin{proposition}[Gap in the specified projector normalization]
For a connected closed oriented cellulation containing at least one edge
with distinct endpoints, the smallest positive eigenvalue of $H_{\rm code}$
is two. In particular, this applies to all four simulated torus and
connected-sum cellulations in this chapter.
\end{proposition}
\begin{proof}
Commuting Hermitian involutions admit a simultaneous eigenbasis.  In it,
the energy is the number of violated checks.  The relations
$\prod_vA_v=I$ and $\prod_fB_f=I$ force the number of violations of
each type to be even.  Consequently any nonzero energy is at least two.
Choose an edge with distinct endpoints and apply its $Z$ operator to
$|\Omega\rangle$.  Exactly its two endpoint stars change sign; all face
checks retain eigenvalue $+1$.  This nonzero vector has energy two,
which proves that the lower bound is attained.
\end{proof}

This statement adds operator and normalization data that an Euler
characteristic does not contain.  It is a gap theorem for an explicit
finite commuting-projector $\mathbb Z_2$ gauge model.  It is not the
mass-gap theorem for nonabelian continuum Yang--Mills theory.  Even the
physical questions within the code model require care: a growing error
string can move its endpoints without increasing the number of defects.
A gap alone therefore does not prove a passive, thermally self-correcting
memory or a fault-tolerance threshold.  Those conclusions need dynamical
noise and recovery assumptions \cite{QSDennis2002}.

\section{A surface-local variational neural circuit}

Cong, Choi and Lukin introduced QCNNs using repeated local unitary layers,
measurement-conditioned pooling, and progressively reduced system size.
Their applications include phase recognition and optimization of quantum
error correction \cite{QSCong2019}.  The present demonstrator is more
specific: it is a trainable four-parameter quantum correlator on a star
of an encoded surface state.  It has executable non-Clifford rotations
and a measured output, but does not reproduce that paper's multiscale
QCNN, learned decoder, or claimed performance comparisons.

Let the first star in the eighteen-qubit cellulation have ordered support
\[
 (q_0,q_1,q_2,q_3)=(0,2,9,15).
\]
The classical input is $s=(s_0,s_1,s_2,s_3)\in\{-1,1\}^4$, and the
trainable vector is $\theta\in\mathbb R^4$.  Define
\[
 \phi_i=s_i\pi/2+\theta_i,\qquad
 R_Y(\phi_i)=\exp(-i\phi_iY/2).
\]
Apply these rotations to the four listed qubits of $|\Omega\rangle$.
Then apply CNOTs with controls $q_1,q_2,q_3$ and common target $q_0$.
The output is the expectation of $Z_{q_0}$.  Four-edge stars make this
circuit local in the cellulation's edge-incidence sense; a physical chip
still needs a specified coupling graph to realize those interactions.

\begin{theorem}[Exact output of the variational correlator]
For the specified state, patch and gates,
\begin{equation}
 f_\theta(s)=\left(\prod_{i=0}^3s_i\right)
              \left(\prod_{i=0}^3\cos\theta_i\right).
 \label{eq:qnn-exact-output}
\end{equation}
\end{theorem}
\begin{proof}
Conjugating $Z_{q_0}$ backward through the three CNOTs gives
$Z_{q_0}Z_{q_1}Z_{q_2}Z_{q_3}$.  Direct multiplication of the
two-by-two rotation matrix gives
\[
 R_Y(\phi)^\dagger ZR_Y(\phi)=\cos\phi\,Z-\sin\phi\,X.
\]
Expand the product of these four expressions.  Any term other than the
all-$X$ term contains a nonempty subset of $Z$ factors on the star.
The four neighbors of its center are distinct.  A chosen $Z$ edge is
incident to a neighboring vertex with no other edge in that subset.
Thus its $Z$ mask has nonzero incidence and is not in $\ker H_X$.

More explicitly, for a Pauli $X(a)Z(b)$, the expansion of
$\langle\Omega|X(a)Z(b)|\Omega\rangle$ is zero unless $a\in W$;
when $a\in W$, it is $|W|^{-1}\sum_{w\in W}(-1)^{b\cdot w}$.
If $b$ is not orthogonal to $W$, pairing $w$ with $w+w_0$ for a
vector satisfying $b\cdot w_0=1$ makes this sum zero.  Therefore all
the terms just identified have zero expectation.

The all-$X$ product is exactly the star stabilizer and has expectation
one.  Its coefficient is $\prod_i(-\sin\phi_i)=\prod_i\sin\phi_i$,
because there are four factors.  Finally
$\sin(s_i\pi/2+\theta_i)=s_i\cos\theta_i$, proving the formula.
\end{proof}

This theorem makes the neural experiment independently checkable.  A
state-vector implementation must execute the rotations and CNOTs, and
its result can be compared with the product formula.  The formula itself
is an efficient classical baseline.  Training this circuit therefore
does not demonstrate a quantum computational advantage or solve a
previously unknown classification problem.

\section{Training, parameter shifts and the meaning of a held-out input}

The target is the four-bit parity $y(s)=\prod_i s_i$.  For a nonempty
training set $T$, use mean squared error
\[
 \mathcal L(\theta)=|T|^{-1}\sum_{s\in T}(f_\theta(s)-y(s))^2.
\]
Let $C(\theta)=\prod_i\cos\theta_i$.  The exact output theorem gives
\begin{equation}
 \mathcal L(\theta)=(C(\theta)-1)^2,
 \qquad
 \partial_j\mathcal L
 =-2(C-1)\sin\theta_j\prod_{i\ne j}\cos\theta_i.
 \label{eq:qnn-loss-gradient}
\end{equation}
These identities include points where $\cos\theta_j=0$, unlike an
expression written using $\tan\theta_j$ and division by that cosine.
They supply an independent reference gradient.

The quantum parameter-shift identity for this circuit is
\[
 \partial_j f_\theta(s)
 =\tfrac12\bigl(f_{\theta+(\pi/2)e_j}(s)
                -f_{\theta-(\pi/2)e_j}(s)\bigr).
\]
To prove it directly, hold the other parameters fixed in
\eqref{eq:qnn-exact-output} and apply
$[\cos(t+\pi/2)-\cos(t-\pi/2)]/2=-\sin t$.
More generally this two-evaluation rule follows when a parameter appears
once in a gate generated by a Pauli operator divided by two.  Repeated
occurrences require a sum of gatewise derivative contributions, and
different generator spectra require an appropriate rule
\cite{QSSchuld2019}.  The shift formula applies to the expectation;
the loss derivative then uses the ordinary chain rule.  Shifting the
entire squared loss by $\pm\pi/2$ is generally not the same calculation.

The recorded training specification uses initial angles
$(0.65,-0.4,0.5,-0.6)$, four sign inputs
$++++,-+++,--++,+-++$, twenty-four gradient steps, and learning rate
$0.2$.  All sixteen possible sign inputs are then evaluated, including
twelve not used in the numerical updates.  The theorem explains why the
same learned scale controls every input: the loss is independent of which
nonempty subset of the truth table was selected.  Consequently the
held-out results are useful software checks, but they cannot support a
statistical generalization claim for an unknown data distribution.

If the output were measured physically, each shot would produce $+1$ or
$-1$, with mean $f_\theta(s)$.  For $M$ independent shots, the sample
mean has variance $(1-f_\theta(s)^2)/M$.  A full-state simulator can
compute the expectation directly without sampling noise.  Its gradient
agreement therefore verifies numerical circuit implementation, not the
sample complexity or robustness of training on a noisy device.  The
report distinguishes these modes explicitly.

\section{Execution evidence and scaling}

The CPU validation executes the forty-two encoder gates on a complex128
array of length $2^{18}=262144$.  It compares the result with an
independently constructed row-space superposition and evaluates all nine
star and nine face expectations.  The recorded norm squared is
$0.9999999999999989$, the maximum amplitude discrepancy is below
$4.2\times10^{-17}$, and the state has exactly 256 nonzero amplitudes
at the expected locations.  The sixteen variational input computations
agree with \eqref{eq:qnn-exact-output} to below $1.1\times10^{-15}$.
These are floating-point execution results; the encoder and output
theorems above provide the exact mathematical assertions.

The topology checks are independent of floating-point geometry.  They
verify edge incidences, oriented face closure, connectivity, rotation
cycles, Euler characteristic, commutation and binary matrix ranks.  The
genus-two and genus-three certificates therefore establish nontrivial
CSS codes even when a browser displays their charts without storing
their exponentially large state vectors.  A visual drawing can be
changed without changing the checked combinatorial data.

The sparse grid interface constructs star and face supports in
$O(L_xL_y)$ storage.  It can supply much larger toric instances to a
stabilizer backend without allocating a dense parity-check matrix or
performing generic binary elimination.  The rank formula is then a
proved topological result, and the manifest identifies whether ranks
were computed by row reduction or obtained from that theorem.  The
2048-qubit archived case additionally has exact row-reduction checks.

A complex64 state vector requires $8\cdot2^n$ bytes before workspaces.
Thirty qubits require 8 GiB; thirty-one require 16 GiB.  A measurement
ancilla can therefore double the dominant allocation.  Forty-six
qubits would require 512 TiB for the raw state alone.  These arithmetic
facts explain why different backends are necessary.  A compressed
Clifford state or batch of Pauli errors can be processed at a much
larger qubit count, but arbitrary $R_Y$ rotations generally leave the
Clifford class.

\subsection{Recorded GPU execution}

\IfFileExists{chapters/quantum-gpu-results.tex}{
 The executable GPU evidence was obtained on 13 November 2026 on an
NVIDIA GeForce RTX 4090.  The recorded environment reports
25,250,627,584 bytes of device memory and 23,960,354,816 initially free
bytes.  It uses Python 3.12.3, cuQuantum Python 26.11.0, cuStateVec
1.15.1, cuStabilizer 0.5.0, and CuPy 14.2.0.  The original record is
\path{extension/quantum/results/gpu-results.json}; it contains nine
successful experiments, the numerical outputs, timestamps, package
versions, and the SHA256 of the executed script.  The source and the
record are retained separately so that a later, larger test does not
silently replace earlier evidence.

\paragraph{Dense state-vector results.}
A two-qubit Bell-state calculation first checked amplitude ordering and
$XX,ZZ$ expectations.  The surface-code runs then executed the emitted
Hadamard/CNOT encoder using actual cuStateVec matrix applications.  They
computed all star and face expectations, injected an $X$ error on a
specified edge, and compared every resulting expectation with the
predicted incident-face signs.  One thirty-qubit run appended a fresh
ancilla and four data-to-ancilla CNOTs to extract one face syndrome.
This final experiment explicitly stored a thirty-one-qubit state.

\begin{center}
\begin{tabular}{lrrrr}
\hline
Data and ancillas & Raw state & Encoder gates & Preparation (s) & Timed total (s)\\
\hline
$18+0$ & 2 MiB & 42 & 0.000312 & 0.000537\\
$30+0$ & 8 GiB & 74 & 0.918804 & 1.204211\\
$30+1$ & 16 GiB & 74 & 1.838463 & 2.528538\\
\hline
\end{tabular}
\end{center}

All three cases use complex64 amplitudes.  Their recorded norm squared
is approximately $0.9999997616$ and their maximum pre-error stabilizer
expectation discrepancy is $2.385\times10^{-7}$.  The extra ancilla has
$Z$ expectation approximately $-0.9999997616$, as predicted by the
injected error.  The routine evaluates these expectations directly; it
does not simulate a finite collection of measurement shots.  The listed
raw-state sizes are exact allocations implied by the precision and
qubit count.  The required gate workspace was zero in these runs.

Preparation time covers encoder application and synchronization.  The
timed total begins after construction of the state-vector object and
includes preparation, the reported expectation checks, and the inserted
error and extraction gates.  It excludes state allocation and final
cleanup.  These are measured intervals for this particular implementation,
not complete application launch times or portable hardware benchmarks.

\paragraph{Training a circuit that has an exact analytic baseline.}
The eighteen-qubit neural experiment executed the encoded preparation,
rotations, fan-in CNOTs, and output expectation on the GPU.  Each gradient
component was obtained from the parameter-shift expectation evaluations;
the analytic product formula was used as a validation reference, not as
a replacement for the quantum gates.  Twenty-four updates at learning
rate $0.2$ changed the initial mean squared error
$0.21987940665826008$ to $0.0027618471230164243$, a reduction by a
factor of approximately $79.6$.  The recorded final angles are
\[
 (0.1918860335,-0.1237818552,0.1522945504,-0.1791698376).
\]
The largest discrepancy from the exact output formula along the recorded
training history was $3.093\times10^{-7}$.  All sixteen final truth-table
entries had the correct sign, including the twelve inputs omitted from
the updates.  The training and evaluation interval was approximately
$0.045774$ seconds, excluding the initial encoded-state construction.

The initial angles already give $C(\theta)>0$.  Consequently the
thresholded sign classifier is correct on all sixteen inputs before
training as well as afterwards.  The measured improvement is the
calibration of expectation magnitude and mean squared error; it is not
an increase from poor classification accuracy to perfect classification.
The magnitude changes from approximately $0.531087$ to $0.947447$.
If a single ideal measurement were used to predict parity, the exact
Born-rule success probability $(1+C)/2$ would accordingly change from
approximately $0.76554$ to $0.97372$.  These probabilities are deductions
from computed expectations, not empirically sampled success frequencies.

\paragraph{Pauli frames and the reference-state assumption.}
For the larger experiments, cuStabilizer stores batched Pauli errors
relative to a known ideal encoded reference.  It does not store the
$2^n$ amplitudes of that reference.  The mathematical encoder theorem
supplies an ideal state with every specified check equal to $+1$.
The frame experiments assume that state and propagate its errors through
ideal syndrome extraction.  They do not execute the large encoder or
construct a generic stabilizer tableau for an arbitrary unknown state.

Each experiment uses 32,768 error realizations with seed 20261113.
Every data qubit receives an independent $X$ error with probability
$0.001$ and an independent $Z$ error with the same probability.  Both
can occur, giving the corresponding $Y$ error up to an irrelevant phase.
Syndrome ancillas and extraction gates are noiseless.  A star ancilla
is reset, Hadamard transformed, used as the control for its incident
CNOTs, Hadamard transformed again, and measured.  A face ancilla receives
CNOTs controlled by the boundary data qubits and is then measured.
Measurement randomization is disabled because the stored output is the
error-induced sign change from the assumed deterministic ideal outcome.

Let $x,z\in\mathbb F_2^E$ be the data error masks.  The exact predicted
syndrome is
\[
 s_X=H_Xz,\qquad s_Z=H_Zx.
\]
With these noiseless, freshly reset ancillas, the extraction circuit
preserves the data error masks.  Hence the implementation can compare
its measured sign-change bits with sparse binary products computed from
the returned data frames.  It does so for every stored syndrome byte,
not merely for a random sample.  Every comparison in the recorded runs
has zero mismatches.  This checks the consistency of the extraction
circuit, qubit indexing, packed-bit representation and binary incidence
calculation.  It does not test a recovery algorithm or estimate a logical
failure probability.

\begin{center}
\begin{tabular}{lrrrr}
\hline
Case & Data & Ancillas & Apply interval (s) & Timed total (s)\\
\hline
Genus two & 32 & 30 & 0.004756 & 3.025362\\
Torus $L=32$ & 2,048 & 2,048 & 0.005072 & 0.013088\\
Torus $L=128$ & 32,768 & 32,768 & 0.073878 & 0.165669\\
Torus $L=512$ & 524,288 & 524,288 & 1.147181 & 2.693085\\
\hline
\end{tabular}
\end{center}

Here the apply interval is a synchronized wall-clock interval around
noise and extraction circuit application, including host dispatch.  It
is not a CUDA-event measurement of kernel time alone.  The timed total
includes simulator setup and exhaustive syndrome verification, but
excludes topology generation before entry into the timer and final
cleanup afterwards.  First-use compilation and initialization make the
genus-two total much larger than its apply interval.  No cross-backend
speedup is inferred from these differently scoped measurements.

For the $L=512$ experiment, all 2,147,483,648 packed syndrome bytes were
checked.  There were 68,489,748 violated checks over all shots, giving
an aggregate syndrome rate of approximately $0.00398663$.  For a
weight-four check under independent relevant Pauli errors of probability
$p$, the exact odd-parity probability is
\[
 \frac{1-(1-2p)^4}{2}.
\]
At $p=0.001$ this is approximately $0.00398802$.  This agreement is a
useful noise-model sanity check.  Correlations among neighboring checks
and their global parity constraints mean that the aggregate count should
not be interpreted as independent Bernoulli samples at every check.

\paragraph{Memory accounting and the scope of a capacity claim.}
For $S$ shots, the stored row stride in bytes is
$4\lceil S/32\rceil$.  With $N$ total data-plus-ancilla qubits and $A$
measured ancillas, the implementation estimates the explicit packed bit
tables as
\[
 B=(2N+A)\,4\lceil S/32\rceil.
\]
The factor two accounts for the $X$ and $Z$ frame components.  This is
an analytic allocation estimate, not an instrumented peak-memory
measurement.  The first run's $L=512$ case has $B=10$ GiB.  Before
allocation, the program requires $1.25B$ to fit within available device
memory after retaining a fixed 2 GiB reserve.  This policy leaves room
for temporary arrays and the desktop; it does not establish a universal
memory bound for every library configuration.

A separate followup script, \path{extension/quantum/gpu_capacity.py},
retains the original experiment source and result unchanged.  It hashes
its own source, the imported experiment source, the topology source,
all archived case JSON files, and the original result.  Its capacity
rule fixes 32,768 shots and chooses the largest multiple of 32 satisfying
\[
 1.25\,(10L^2)(4096)\leq
       \text{free device bytes}-2\,\mathrm{GiB}.
\]
The resulting test is a bounded workload selected under stated memory
assumptions.  A different shot count, error representation, precision,
number of observables, or available-memory state changes the achievable
qubit count.  The dense thirty-one-qubit result and the much larger
frame results therefore answer different computational questions.

\paragraph{Separate capacity followup.}
The followup ran from 10:58:12 to 10:58:17 UTC on the same date and
device.  It successfully exercised the genus-three code with 46 data
qubits and 42 syndrome ancillas, and then selected $L=640$ from the
observed free-memory value of 23,922,606,080 bytes.  The latter case has
819,200 data qubits and 819,200 ancillas, for a total of 1,638,400
qubits in each of 32,768 Pauli frames.  Its estimated packed tables
occupy 16,777,216,000 bytes, or 15.625 GiB.  The budget including the
25 percent overhead allowance and 2 GiB reserve is 23,119,003,648
bytes, below the recorded available amount.

\begin{center}
\begin{tabular}{lrrrr}
\hline
Case & Apply (s) & Timed body (s) & Full call (s) & Checked bytes\\
\hline
Genus three & 0.002198 & 0.081878 & 0.209550 & 172,032\\
Torus $L=640$ & 1.791833 & 4.327947 & 4.813268 & 3,355,443,200\\
\hline
\end{tabular}
\end{center}

Both cases had zero syndrome-byte mismatches.  The larger test observed
107,029,374 violated checks, an aggregate rate of approximately
$0.00398715$.  The additional full-call timer includes topology
generation and cleanup; it still excludes Python startup, library
imports, and output-file serialization.  These boundaries make the
measurements reproducible without calling a host wall-clock interval
a pure kernel time.

The separate evidence file is
\path{extension/quantum/results/capacity-results.json}.  Its source and
case hashes agree with the archived local inputs, and it records that
those inputs were unchanged during execution.  A byte-exact copy of the
original first-run JSON is retained as
\path{extension/quantum/results/gpu-results-remote-original.json}; the
ordinary first-run JSON has identical parsed values but a few different
numeric formatting choices.  This distinction explains their different
file hashes without conflating it with a difference in experiment data.
The largest completed frame workload is therefore the stated $L=640$
case under the stated memory policy.  This observation establishes
neither an absolute maximum qubit count nor the feasibility of arbitrary
non-Clifford computation at that scale.

}{
 The GPU execution report is generated separately from this mathematical
 chapter.  Until that report is included, no GPU timing, peak-memory
 measurement or completed large-instance run is asserted here.  The
 executable cases, gate conventions and CPU cross-checks above specify
 what an additional backend must reproduce.
}

\section{Extensions with explicit mathematical obligations}

For a prime $p$, the cellular construction extends to qudits of dimension
$p$.  Retain signed integer incidence matrices and reduce them modulo
$p$.  Define $X|j\rangle=|j+1\rangle$ and
$Z|j\rangle=e^{2\pi i j/p}|j\rangle$.  Their products commute exactly
when the corresponding signed row pairing is zero modulo $p$.
The integer identity $\partial_1\partial_2=0$ supplies this condition,
and the connected oriented-surface rank argument again gives $2g$
logical qudits.  Over composite dimensions, field rank arguments must
be replaced by module or Smith-normal-form calculations.  This is an
actual point where integral order and module theory helps prevent an
invalid inference.  The present executable gate engine remains a qubit
implementation; the qudit construction is a mathematical extension.

Nonabelian finite-group gauge models require richer local Hilbert spaces
and constraints.  Kitaev's anyon construction develops that direction,
but the qubit star-and-face specification cannot simply be relabelled as
a universal nonabelian simulator \cite{QSKitaev2003}.  Likewise an
arbitrary hypermap encoding may admit a circuit transformation to a
surface code, rather than introducing a fundamentally different source
of protection; constructive comparisons are available in the literature
\cite{QSSarvepalli2014}.

For learned quantum circuits, a substantive next step would specify a
noise channel, a family of input states, a training distribution, a
decoder or observable, and comparison with classical and analytic
baselines.  Multi-layer pooling and measurement feedback would require
additional circuit semantics beyond the present four-parameter model.
Any claim of improved decoding should measure logical failure after
recovery, not just a reduced training loss or an attractive surface
visualization.

The connection to the rest of the monograph is now concrete.  A clean
dessin supplies a rotation system and an integral cell complex.  Reduction
modulo two supplies compatible binary matrices.  Their zero composite
becomes commutation of quantum operators.  Homology becomes logical
Pauli operators.  A row-space construction becomes a finite gate list,
and a local observable becomes a trainable circuit with a provable output.
Every transition has specified data and a checkable statement.  These
are useful new executable interfaces within the project, without
identifying them with a solution of GRH, an improved ECC attack, or a
continuum Yang--Mills mass-gap proof.
